% VI 文本编辑器
% 使用xelatex编译

\documentclass[12pt,a4paper,twoside]{ctexbook}

% 页面设置
\input{../../../../../Book/0_book/common/page}

% 字体设置
\usepackage{xeCJK}
\usepackage{fontspec}
\usepackage{microtype}

% 设置中文字体
\setCJKmainfont{SimSun}[
    BoldFont=SimHei,
    ItalicFont=KaiTi
]
\setCJKsansfont{SimHei}
\setCJKmonofont{SimSun}
\setCJKfamilyfont{kai}[
    BoldFont=KaiTi
]{KaiTi}
\setCJKfamilyfont{fs}[
    BoldFont=FangSong
]{FangSong}

% 常用字体命令
\newcommand{\song}{\CJKfamily{zhsong}}
\newcommand{\hei}{\CJKfamily{zhhei}}
\newcommand{\kai}{\CJKfamily{kai}}
\newcommand{\fs}{\CJKfamily{fs}}

% 标题格式设置
\ctexset{
    part/name={第,卷},
    part/number={\Roman{part}},
    chapter/name={第,章},
    chapter/number={\chinese{chapter}},
    section/name={第,节},
    section/number={\arabic{section}},
    subsection/number={\arabic{section}.\arabic{subsection}},
    chapter/format={\centering\hei\zihao{2}},
    section/format={\hei\zihao{4}},
    subsection/format={\hei\zihao{5}}
}

% 目录设置
\usepackage{titletoc}
\titlecontents{chapter}[0pt]{\vspace{10pt}\bfseries\zihao{-4}}{\contentspush{\thecontentslabel\hspace{1em}}}{}{\titlerule*[8pt]{.}\contentspage}
\titlecontents{section}[2.5em]{\vspace{5pt}\zihao{5}}{\contentspush{\thecontentslabel\hspace{1em}}}{}{\titlerule*[8pt]{.}\contentspage}
\titlecontents{subsection}[5em]{\vspace{3pt}\zihao{5}}{\contentspush{\thecontentslabel\hspace{1em}}}{}{\titlerule*[8pt]{.}\contentspage}

% 页眉页脚设置
\usepackage{fancyhdr}
\pagestyle{fancy}
\fancyhf{}
\fancyhead[LE,RO]{\zihao{5}\thepage}
\fancyhead[LO]{\zihao{5}\leftmark}
\fancyhead[RE]{\zihao{5}\rightmark}
\renewcommand{\chaptermark}[1]{\markboth{\chaptername\ \thechapter\ #1}{}}
\renewcommand{\sectionmark}[1]{\markright{\thesection\ #1}}
\fancyfoot[C]{\zihao{5} \thepage}
\renewcommand{\headrulewidth}{0.4pt}
\renewcommand{\footrulewidth}{0pt}

% 插图设置
\usepackage{graphicx}
\usepackage{float}
\usepackage{subfigure}
\graphicspath{{images/}}
\floatstyle{plaintop}
\restylefloat{figure}

% 表格设置
\usepackage{tabularx}
\usepackage{booktabs}
\usepackage{longtable}

% 数学公式设置
\usepackage{amsmath, amssymb, amsthm}
\usepackage{mathrsfs}

% 定理环境
\newtheorem{theorem}{定理}[chapter]
\newtheorem{definition}{定义}[chapter]
\newtheorem{lemma}{引理}[chapter]
\newtheorem{corollary}{推论}[chapter]
\newtheorem{example}{例}[chapter]

% 引用设置
\usepackage{hyperref}
\hypersetup{
    colorlinks=true,
    linkcolor=blue,
    citecolor=blue,
    urlcolor=blue,
    pdftitle={VI 文本编辑器},
    pdfauthor={Rocky Linux},
    pdfsubject={VI 文本编辑器},
    pdfkeywords={VI, VIM, 文本编辑器, Linux}
}

% 目录深度
\setcounter{tocdepth}{3}
\setcounter{secnumdepth}{3}

% 自定义环境
\newenvironment{note}{
    \begin{center}
    \fbox{\parbox{0.9\linewidth}{
        \vspace{0.5em}
        \textbf{注意：}
        \vspace{0.5em}
    }}
    \end{center}
}{
}

\newenvironment{tip}{
    \begin{center}
    \fbox{\parbox{0.9\linewidth}{
        \vspace{0.5em}
        \textbf{提示：}
        \vspace{0.5em}
    }}
    \end{center}
}{
}

\newenvironment{warning}{
    \begin{center}
    \fbox{\parbox{0.9\linewidth}{
        \vspace{0.5em}
        \textbf{警告：}
        \vspace{0.5em}
    }}
    \end{center}
}{
}

\begin{document}

\title{\hei\zihao{0} VI 文本编辑器}
\author{\song\zihao{2} Rocky Linux}
\date{\song\zihao{4} \today}

\begin{titlepage}
    \begin{center}
        \vspace*{6cm}
        \hei\zihao{0} VI 文本编辑器
        \vspace*{3cm}
        \song\zihao{2} Rocky Linux
        \vspace*{3cm}
        \song\zihao{4} \today
    \end{center}
\end{titlepage}

\tableofcontents

\mainmatter

\chapter{VI 文本编辑器}

在本章中，您将学习如何使用VI（Visual）编辑器。

\section{学习目标}

在本章中，未来的Linux管理员将学习如何：

\begin{itemize}
    \item 使用VI编辑器的主要命令
    \item 使用VI编辑器修改文本
\end{itemize}

\textbf{关键词：} 用户命令，linux

\textbf{知识要求：} \textstar\ \textstar

\textbf{复杂度：} \textstar\ \textstar\ \textstar

\textbf{阅读时间：} 20分钟

\section{简介}

\textit{Visual}（\textit{VI}）是Linux下非常流行的文本编辑器，尽管其人机工程学有限。它确实是一个完全基于文本模式的编辑器：每个操作都通过键盘上的按键或专用命令完成。

它非常强大，最重要的是非常实用，因为它对于基本应用程序来说是最小化的。因此在系统故障时也可以访问它。它的\textit{通用性}（它存在于所有Linux发行版和Unix下）使其成为管理员的\textit{关键}工具。

它的功能包括：

\begin{itemize}
    \item 插入、删除、修改文本
    \item 复制单词、行或文本块
    \item 搜索和替换字符
\end{itemize}

\section{vi命令}

\texttt{vi}命令打开\textit{VI}文本编辑器。

\begin{verbatim}
vi [-c command] [file]
\end{verbatim}

示例：

\begin{verbatim}
$ vi /home/rockstar/file
\end{verbatim}

\begin{table}[H]
    \centering
    \caption{vi命令选项}
    \begin{tabular}{ll}
        \toprule
        \textbf{选项} & \textbf{说明} \\
        \midrule
        \texttt{-c command} & 在打开时通过指定命令执行VI \\
        \bottomrule
    \end{tabular}
\end{table}

如果文件存在于路径指定的位置，它将被VI读取，VI将处于\textbf{命令}模式。

如果文件不存在，VI将打开一个空白文件，屏幕上显示一个空白页。当文件保存时，它将采用命令中指定的名称。

如果执行\texttt{vi}命令时未指定文件名，VI将打开一个空白文件，屏幕上显示一个空白页。当文件保存时，VI将要求提供文件名。

\texttt{vim}编辑器采用VI的界面和功能，并进行了许多改进。

\begin{verbatim}
vim [-c command] [file]
\end{verbatim}

在这些改进中，用户拥有语法高亮，这对于编辑shell脚本或配置文件非常有用。

vim是Visual Interface的简称，它可以执行输出、删除、查找、替换、块操作等文本操作，而且用户可以根据自己的需要对其进行定制。这是其他编辑程序没有的。vim不是一个排版程序，不可以对字体、格式、段落等其他属性进行编排，只是一个文本编辑程序。vim是全屏幕文本编辑器，没有菜单，只有命令。

在会话期间，VI使用一个缓冲区文件，在其中记录用户所做的所有更改。

\begin{note}
只要用户尚未保存其工作，原始文件就不会被修改。
\end{note}

启动时，VI处于\textit{命令}模式。

\begin{tip}
文本行通过按ENTER键结束，但如果屏幕不够宽，VI会自动换行，默认为\textit{wrap}配置。这些换行可能不是想要的，这就是\textit{nowrap}配置。
\end{tip}

要退出VI，从命令模式，按\texttt{:}然后输入：

\begin{itemize}
    \item \texttt{q} 退出而不保存（\textit{quit}）
    \item \texttt{w} 保存您的工作（\textit{write}）
    \item \texttt{wq}（\textit{write quit}）或\texttt{x}（\textit{eXit}）保存并退出
\end{itemize}

要强制退出而不确认，必须在上述命令中添加\texttt{!}。

\begin{warning}
没有定期备份，因此您必须记得定期保存您的工作。
\end{warning}

\section{操作模式}

在VI中，有3种工作模式：

\begin{itemize}
    \item \textit{命令}模式
    \item \textit{插入}模式
    \item \textit{ex}模式
\end{itemize}

VI的理念是在\textit{命令}模式和\textit{插入}模式之间交替。

第三种模式\textit{ex}是来自旧文本编辑器的底部命令模式。

\subsection{命令模式}

这是VI启动时的默认模式。要从任何其他模式访问它，只需按ESC键。

所有输入都被解释为命令，并执行相应的操作。这些主要是编辑文本的命令（复制、粘贴、撤销等）。

命令不会显示在屏幕上。

\subsection{插入模式}

这是文本修改模式。要从\textit{命令}模式访问它，必须按下将执行操作并更改模式的特殊键。

文本不会直接输入到文件中，而是输入到内存中的缓冲区。只有在保存文件时，更改才会生效。

\subsection{Ex模式}

这是文件修改模式。要访问它，必须先切换到\textit{命令}模式，然后输入通常以字符\texttt{:}开头的\textit{ex}命令。

通过按ENTER键验证命令。

\section{移动光标}

在\textit{命令}模式下，有几种移动光标的方法。

鼠标在文本环境中不活动，但在图形环境中是活动的，可以逐个字符移动它，但存在快捷键可以更快地移动。

移动光标后，VI保持在\textit{命令}模式。

光标放置在所需字符下方。

\subsection{从字符移动}

\begin{itemize}
    \item 向左移动一个或\texttt{n}个字符：\texttt{$\leftarrow$}、\texttt{n$\leftarrow$}、\texttt{h}或\texttt{nh}
    \item 向右移动一个或\texttt{n}个字符：\texttt{$\rightarrow$}、\texttt{n$\rightarrow$}、\texttt{l}或\texttt{nl}
    \item 向上移动一个或\texttt{n}个字符：\texttt{$\uparrow$}、\texttt{n$\uparrow$}、\texttt{k}或\texttt{nk}
    \item 向下移动一个或\texttt{n}个字符：\texttt{$\downarrow$}、\texttt{n$\downarrow$}、\texttt{j}或\texttt{nj}
    \item 移动到行尾：\texttt{\$}或\texttt{END}
    \item 移动到行首：\texttt{0}或\texttt{POS1}
\end{itemize}

\subsection{从单词的首字符移动}

单词由字母或数字组成。标点字符和撇号分隔单词。

如果光标在单词中间，\texttt{w}移动到下一个单词，\texttt{b}移动到单词的开头。

如果行结束，VI会自动转到下一行。

\begin{itemize}
    \item 向右移动一个或\texttt{n}个单词：\texttt{w}或\texttt{nw}
    \item 向左移动一个或\texttt{n}个单词：\texttt{b}或\texttt{nb}
\end{itemize}

\subsection{从行上的任意位置移动}

\begin{itemize}
    \item 移动到文本的最后一行：\texttt{G}
    \item 移动到第\texttt{n}行：\texttt{nG}
    \item 移动到屏幕的第一行：\texttt{H}
    \item 移动到屏幕的中间行：\texttt{M}
    \item 移动到屏幕的最后一行：\texttt{L}
\end{itemize}

\section{插入文本}

在\textit{命令}模式下，有几种插入文本的方法。

输入这些键之一后，VI切换到\textit{插入}模式。

\begin{note}
VI切换到\textit{插入}模式。因此，您必须按ESC键返回\textit{命令}模式。
\end{note}

\subsection{相对于字符插入}

\begin{itemize}
    \item 在字符之前插入文本：\texttt{i}（\textit{insert}）
    \item 在字符之后插入文本：\texttt{a}（\textit{append}）
\end{itemize}

\subsection{相对于行插入}

\begin{itemize}
    \item 在行首插入文本：\texttt{I}
    \item 在行尾插入文本：\texttt{A}
\end{itemize}

\subsection{相对于文本插入}

\begin{itemize}
    \item 在行之前插入文本：\texttt{O}
    \item 在行之后插入文本：\texttt{o}
\end{itemize}

\section{字符、单词和行}

VI通过管理以下内容来编辑文本：

\begin{itemize}
    \item 字符
    \item 单词
    \item 行
\end{itemize}

在每种情况下，都可以：

\begin{itemize}
    \item 删除
    \item 替换
    \item 复制
    \item 剪切
    \item 粘贴
\end{itemize}

这些操作在\textit{命令}模式下完成。

\subsection{字符}

\begin{itemize}
    \item 删除一个或\texttt{n}个字符：\texttt{x}或\texttt{nx}
    \item 用另一个字符替换一个字符：\texttt{r character}
    \item 用其他字符替换多个字符：\texttt{R characters ESC}
\end{itemize}

\begin{note}
R命令切换到\textit{替换}模式，这是一种\textit{插入}模式。
\end{note}

\subsection{单词}

\begin{itemize}
    \item 删除（剪切）一个或\texttt{n}个单词：\texttt{dw}或\texttt{ndw}
    \item 复制一个或\texttt{n}个单词：\texttt{yw}或\texttt{nyw}
    \item 在光标后粘贴一个单词一次或\texttt{n}次：\texttt{p}或\texttt{np}
    \item 在光标前粘贴一个单词一次或\texttt{n}次：\texttt{P}或\texttt{nP}
    \item 替换一个单词：\texttt{cw word ESC}
\end{itemize}

\begin{tip}
必须将光标定位在要剪切（或复制）的单词的第一个字符下，否则VI将只剪切（或复制）光标和单词末尾之间的单词部分。删除单词就是剪切它。如果之后不粘贴，缓冲区将被清空，单词将被删除。
\end{tip}

\subsection{行}

\begin{itemize}
    \item 删除（剪切）一个或\texttt{n}行：\texttt{dd}或\texttt{ndd}
    \item 复制一个或\texttt{n}行：\texttt{yy}或\texttt{nyy}
    \item 在当前行之后粘贴已复制或删除的内容一次或\texttt{n}次：\texttt{p}或\texttt{np}
    \item 在当前行之前粘贴已复制或删除的内容一次或\texttt{n}次：\texttt{P}或\texttt{nP}
    \item 从行首删除（剪切）到光标：\texttt{d0}
    \item 从光标删除（剪切）到行尾：\texttt{d\$}
    \item 从行首复制到光标：\texttt{y0}
    \item 从光标复制到行尾：\texttt{y\$}
    \item 从当前行删除（剪切）文本：\texttt{dL}或\texttt{dG}
    \item 从当前行复制文本：\texttt{yL}或\texttt{yG}
\end{itemize}

\subsection{取消操作}

\begin{itemize}
    \item 撤销上一次操作：\texttt{u}
    \item 撤销当前行上的操作：\texttt{U}
\end{itemize}

\subsection{取消撤销}

\begin{itemize}
    \item 取消撤销：\texttt{Ctrl+R}
\end{itemize}

\section{EX命令}

\textit{Ex}模式允许您对文件进行操作（保存、布局、选项等）。搜索和替换命令也在\textit{Ex}模式中输入。命令显示在页面底部，必须按ENTER键验证。

要从\textit{命令}模式切换到\textit{Ex}模式，输入\texttt{:}。

\subsection{行号显示}

\begin{itemize}
    \item 显示/隐藏行号：\texttt{:set nu}和更长的\texttt{:set number}
    \item 隐藏行号：\texttt{:set nonu}和更长的\texttt{:set nonumber}
\end{itemize}

\subsection{搜索字符串}

\begin{itemize}
    \item 从光标搜索字符串：
    \begin{verbatim}
/string
    \end{verbatim}
    \item 从光标之前搜索字符串：
    \begin{verbatim}
?string
    \end{verbatim}
    \item 转到下一个找到的匹配项：\texttt{n}
    \item 转到上一个找到的匹配项：\texttt{N}
\end{itemize}

VI中有通配符可以方便搜索。

\begin{itemize}
    \item \texttt{[]}：搜索字符范围或单个字符，其可能的值已指定。
    
    示例：
    \begin{itemize}
        \item \texttt{/[Ww]ord}：搜索\textit{word}和\textit{Word}
        \item \texttt{/[1-9]word}：搜索\textit{1word}、\textit{2word} ... \textit{xword}，其中\texttt{x}是一个数字
    \end{itemize}
    
    \item \texttt{\^{}}：搜索以该字符串开头的行。
    
    示例：
    \begin{verbatim}
/^Word
    \end{verbatim}
    
    \item \texttt{\$}：搜索以该字符串结尾的行。
    
    示例：
    \begin{verbatim}
/Word$
    \end{verbatim}
    
    \item \texttt{.}：搜索具有未知字母的单词。
    
    示例：
    \begin{itemize}
        \item \texttt{/W.rd}：搜索\textit{Word}、\textit{Ward}等
    \end{itemize}
    
    \item \texttt{*}：搜索一个或多个字符，无论它们是什么。
    
    示例：
    \begin{verbatim}
/W*d
    \end{verbatim}
\end{itemize}

\subsection{替换字符串}

从文本的第一行到最后一行，将搜索的字符串替换为指定的字符串：

\begin{verbatim}
:1,$ s/search/replace
\end{verbatim}

\textbf{注意：}您也可以使用\texttt{:0,\$s/search/replace}来指定从文件的绝对开头开始。

从第\texttt{n}行到第\texttt{m}行，将搜索的字符串替换为指定的字符串：

\begin{verbatim}
:n,m s/search/replace
\end{verbatim}

默认情况下，每行只替换找到的第一个匹配项。要强制替换每个匹配项，必须在命令末尾添加\texttt{/g}：

\begin{verbatim}
:n,m s/search/replace/g
\end{verbatim}

浏览整个文件以将搜索的字符串替换为指定的字符串：

\begin{verbatim}
:% s/search/replace
\end{verbatim}

\subsection{文件操作}

\begin{itemize}
    \item 保存文件：
    \begin{verbatim}
:w
    \end{verbatim}
    
    \item 以另一个名称保存：
    \begin{verbatim}
:w file
    \end{verbatim}
    
    \item 将第\texttt{n}行到第\texttt{m}行保存到另一个文件：
    \begin{verbatim}
:n,m w file
    \end{verbatim}
    
    \item 重新加载文件的最后一条记录：
    \begin{verbatim}
:e!
    \end{verbatim}
    
    \item 在光标后粘贴另一个文件的内容：
    \begin{verbatim}
:r file
    \end{verbatim}
    
    \item 退出编辑文件而不保存：
    \begin{verbatim}
:q
    \end{verbatim}
    
    \item 退出在会话期间已修改但未保存的文件：
    \begin{verbatim}
:q!
    \end{verbatim}
    
    \item 退出文件并保存：
    \begin{verbatim}
:wq 或 :x
    \end{verbatim}
\end{itemize}

\section{其他功能}

可以通过指定要为会话加载的选项来执行VI。为此，必须使用\texttt{-c}选项：

\begin{verbatim}
$ vi -c "set nu" /home/rockstar/file
\end{verbatim}

也可以将\textit{Ex}命令输入到放置在用户登录目录中的名为\texttt{.exrc}的文件中。在每次VI或VIM启动时，命令将被读取并应用。

\subsection{vimtutor命令}

有一个学习如何使用VI的教程。可以使用命令\texttt{vimtutor}访问它。

\begin{verbatim}
$ vimtutor
\end{verbatim}

\end{document}
